\chapter{MERN Security}

A MERN stack has attack surface on both ends: the Express API accepts
arbitrary input, and React runs in a browser you don't control. This
chapter covers the threats that matter in practice and the specific
mitigation for each --- most are a single middleware addition.

\section{Cross-Site Scripting (XSS)}

Rendering user-supplied content as raw HTML lets an attacker inject a
\texttt{<script>} that runs in other users' browsers, stealing cookies or
tokens. React escapes text by default, so \texttt{\{task.title\}} is safe;
the risk appears only when you bypass that escaping.

\begin{lstlisting}[language=JavaScript]
// DANGEROUS -- only with fully trusted, sanitized content
<div dangerouslySetInnerHTML={{ __html: task.description }} />
// SAFE -- React escapes this automatically
<div>{task.description}</div>
\end{lstlisting}

\begin{warnbox}[WARNING]
If you must render HTML, sanitize it with \texttt{DOMPurify} first. Never
assume "it's our own users" makes stored content safe to render unescaped.
\end{warnbox}

\section{Cross-Site Request Forgery (CSRF)}

CSRF tricks a logged-in user's browser into submitting a request from
another site using auto-attached cookies --- relevant because JWTs in
HttpOnly cookies (Chapter~19) are sent automatically on every request to
your domain. Mitigate with \texttt{sameSite: "strict"/"lax"} cookies and a
CORS whitelist instead of \texttt{*}.

\begin{lstlisting}[language=JavaScript]
// server.js
import cors from "cors";
app.use(cors({ origin: process.env.CLIENT_URL, credentials: true }));
\end{lstlisting}

\section{NoSQL Injection}

Building a Mongo query straight from a request body allows operator
injection: sending \texttt{\{"\$gt": ""\}} as \texttt{password} can make
\texttt{User.findOne({email, password})} match any document, bypassing
login.

\begin{lstlisting}[language=JavaScript]
// server.js
import mongoSanitize from "express-mongo-sanitize";
app.use(mongoSanitize()); // strips keys starting with "$" or containing "."
\end{lstlisting}

\begin{notebox}[NOTE]
Mongoose schema types reject some bad shapes, but that alone is fragile ---
explicit sanitization middleware closes the gap regardless.
\end{notebox}

\section{Brute Force Login Attempts}

Without a limit, an attacker can script thousands of password guesses per
minute against \texttt{/api/auth/login}. Rate limiting caps attempts per
IP per window.

\begin{lstlisting}[language=JavaScript]
// middleware/rateLimiter.js
import rateLimit from "express-rate-limit";
export const loginLimiter = rateLimit({
  windowMs: 15 * 60 * 1000, max: 10,
  message: { success: false, message: "Too many login attempts. Try again later." },
});
// router.post("/login", loginLimiter, loginUser);
\end{lstlisting}

\begin{tipbox}[TIP]
Apply a strict limiter to \texttt{/login} and \texttt{/register}, and a
looser general limiter elsewhere.
\end{tipbox}

\section{Credential Theft via XSS-Accessible Tokens}

A JWT in \texttt{localStorage} is readable by any JS on the page, including
an injected XSS payload. HttpOnly cookies (Chapter~19) are invisible to
\texttt{document.cookie} and JS entirely, so XSS alone can't steal them.

\section{Insecure Cookies}

An auth cookie without \texttt{secure}/\texttt{httpOnly} can be read by
scripts and intercepted over plain HTTP.

\begin{lstlisting}[language=JavaScript]
res.cookie("token", jwtToken, {
  httpOnly: true,
  secure: process.env.NODE_ENV === "production",
  sameSite: "strict",
  maxAge: 7 * 24 * 60 * 60 * 1000,
});
\end{lstlisting}

\section{Common HTTP Header Hardening}

\texttt{helmet} sets security headers (disabling MIME sniffing, a baseline
CSP) that block several attack classes in one line.

\begin{lstlisting}[language=JavaScript]
import helmet from "helmet";
app.use(helmet());
\end{lstlisting}

\section{Exposed Secrets}

Committing \texttt{.env} (or hardcoding a secret) puts credentials in Git
history forever, even after later deletion. See Chapter~38; the rule is
\texttt{.env} in \texttt{.gitignore} and secrets read only via
\texttt{process.env}.

\section{Malicious File Uploads}

An unrestricted upload endpoint lets an attacker upload disguised
executables or oversized files. Restrict MIME type and size at the multer
layer; never trust the extension alone.

\begin{lstlisting}[language=JavaScript]
// middleware/upload.js
import multer from "multer";
const upload = multer({
  limits: { fileSize: 2 * 1024 * 1024 },
  fileFilter: (req, file, cb) => {
    const allowed = ["image/jpeg", "image/png", "image/webp"];
    if (!allowed.includes(file.mimetype)) return cb(new Error("Only JPEG, PNG, or WEBP images are allowed"));
    cb(null, true);
  },
});
export default upload;
\end{lstlisting}

\section{Input Validation}

Every mitigation above assumes untrusted input is validated at the edge.
See Chapter~23 for the \texttt{express-validator} layer that rejects
malformed requests before a controller or query ever runs.

\section{Summary Table}

\begin{center}
\begin{tabularx}{\textwidth}{p{2.6cm}Xp{2.6cm}}
\toprule
\textbf{Threat} & \textbf{Mitigation} & \textbf{Where} \\
\midrule
XSS & Rely on React's default escaping; sanitize any HTML before \texttt{dangerouslySetInnerHTML} & Ch.~36 (rendering) \\
CSRF & \texttt{sameSite} cookie flag + CORS origin whitelist & Ch.~19, Ch.~37 \\
NoSQL injection & \texttt{express-mongo-sanitize} to strip \texttt{\$}/\texttt{.} keys, plus schema validation & Ch.~23, Ch.~37 \\
Brute force login & \texttt{express-rate-limit} on \texttt{/auth} routes & Ch.~37 \\
Credential theft (XSS reading tokens) & HttpOnly cookies instead of \texttt{localStorage} & Ch.~19 \\
Insecure cookies & \texttt{httpOnly}, \texttt{secure}, \texttt{sameSite} flags & Ch.~19 \\
Header-based attacks & \texttt{helmet()} middleware & Ch.~37 \\
Exposed secrets & \texttt{.env} + \texttt{.gitignore}, never hardcode & Ch.~38 \\
Malicious uploads & MIME/type + size limits in multer \texttt{fileFilter} & Ch.~37 \\
Malformed/invalid input & \texttt{express-validator} at the route layer & Ch.~23 \\
\bottomrule
\end{tabularx}
\end{center}

\whatwhyhow
{A secure MERN API layers helmet headers, rate limiting, input sanitization/validation, HttpOnly+secure+sameSite cookies, CORS restriction, upload restrictions, and secret management, instead of one single defense.}
{Each threat exploits a different gap, so a single fix like "just validate input" leaves the others open.}
{Add \texttt{helmet()} and \texttt{mongoSanitize()} globally, rate-limit auth routes, set cookies with \texttt{httpOnly}/\texttt{secure}/\texttt{sameSite}, whitelist CORS origins, restrict upload MIME types/sizes, and keep secrets in an untracked \texttt{.env}.}
